% 148-17(148쪽) — 등변사다리꼴 ABCD 와 두 대각선 AC, BD 가 이루는 각 $60^\circ$.
% 스캔 화질이 낮아 TikZ 로 다시 그림. 원문과 같이 B 를 왼쪽 아래, C 를 오른쪽 아래, A·D 를 짧은 윗변에 두었다.
% 라벨은 원문에 있는 A·B·C·D·$60^\circ$ 만 둔다 — 넓이와 대각선의 길이(답)는 적지 않는다.
\begin{tikzpicture}[line width=0.5pt, scale=1.02, >=stealth]
  \coordinate (B) at (0,0);
  \coordinate (C) at (2.9,0);
  \coordinate (A) at (0.35,1.47);
  \coordinate (D) at (2.55,1.47);
  \coordinate (I) at (1.45,0.836);
  \draw[line width=0.55pt] (A) -- (B) -- (C) -- (D) -- cycle;
  \draw[line width=0.55pt] (A) -- (C);
  \draw[line width=0.55pt] (B) -- (D);
  \draw[line width=0.4pt] ([shift=(150:0.3)]I) arc (150:210:0.3);
  \node[font=\footnotesize] at ([shift=(180:0.62)]I) {$60^\circ$};
  \node[above left=-3.5pt and -3pt, font=\small] at (A) {$\pt{A}$};
  \node[below left=-3.5pt and -3pt, font=\small] at (B) {$\pt{B}$};
  \node[below right=-3.5pt and -3pt, font=\small] at (C) {$\pt{C}$};
  \node[above right=-3.5pt and -3pt, font=\small] at (D) {$\pt{D}$};
\end{tikzpicture}
